\documentclass[withoutpreface,bwprint]{cumcmthesis}
\geometry{margin=2.5cm}
\usepackage{booktabs,tabularx,url}
\hypersetup{hidelinks}
\setlength{\parindent}{2em}
\title{AI工具使用详情}
\author{}
\date{2026年9月10日}
\begin{document}
\maketitle

\section{AI工具使用说明}

\subsection{所用AI工具名称、版本或型号}

本次竞赛过程中使用的AI工具主要包括月之暗面 kimi K3 和其 cli 工具 kimi code。
kimi主要用于赛题拆解、思路讨论、文字整理和文献检索辅助；code主要用于根据
已经确定的数学模型、约束条件和计算流程辅助生成、修改及检查Python程序。

\subsection{具体使用目的和环节}

AI主要作为建模与程序实现过程中的辅助工具使用，未直接替代参赛人员完成模型选择
和结果判定。

在建模初期，首先由AI辅助拆解赛题四个问题的已知条件、决策变量、约束关系和输出要求，
帮助梳理问题一至问题四之间由确定性调度、预测误差、滚动调整到波动电价的递进关系。
在此基础上，参赛人员结合赛题原文、附件数据和已有数学建模知识确定整体建模路线。

具体数学模型的关键约束、费用计算方式以及主要公式由参赛人员结合题意确定，包括
十分钟尺度的功率--电量换算、供需平衡、储能状态转移、储电量上下界、充放电功率限制、
紧急购电费用、计划购电与调整购电结算等内容。AI主要用于检查这些公式之间的逻辑关系、
协助整理为可编程形式，并根据已经给定的数学表达生成相应程序框架。

在计算阶段，主要使用Code根据参赛人员提供的变量定义、目标函数、约束条件以及
数据处理要求生成或修改Python代码。参赛人员随后检查程序中的时间索引、单位换算、
约束方向、费用倍数和数据读写逻辑，确认后再实际运行程序，对附件数据进行计算和分析。

\subsection{主要提示方式与使用过程}

AI的使用总体遵循``题目拆解--人工建模--代码辅助生成--人工检查--实际运行--结果分析''
的过程，而不是直接要求AI从题目自动生成完整答案。

首先，将赛题原文及各附件含义提供给AI，要求其仅拆解各问的目标、已知信息、
决策时点、费用规则和需要输出的结果，用于帮助理解题意。例如：

\begin{quote}
请先不要直接求解，只分析本题每一问的已知量、决策变量、约束条件、费用组成，
以及四个问题之间的递进关系。
\end{quote}

在明确建模思路后，由参赛人员给出具体数学约束和计算口径，再要求AI将模型转换为程序。
典型提示方式例如：

\begin{quote}
现采用10分钟为一个时段，购电量、充电量和放电量均以kWh表示。
储能状态满足
$E_{t+1}=E_t+0.9C_t-D_t/0.9$，
储电量限制为$1200\leq E_t\leq10800$，
单时段充放电量不超过$5000/6$ kWh。
请按照上述约束生成线性规划代码，不自行增加新的约束和费用项。
\end{quote}

对于问题二至问题四，则继续由参赛人员给出紧急购电、调整购电、滚动更新时间、
历史场景构造等规则，AI根据这些要求补充相应代码。生成程序后，对关键模块逐项检查，
重点核对功率与电量单位、时间索引、储能状态递推、费用计算以及是否误用未来信息。

程序确认后由参赛人员实际运行，并根据运行得到的费用、购电量、储能状态和紧急购电结果
进行比较，再决定是否保留相应策略。对于不同模型或参数组合，主要依据实际计算结果
而不是AI的文字判断进行选择。

\subsection{AI输出的采纳、人工修改和核验情况}

AI输出没有直接作为最终结果使用。对于赛题分析和建模建议，仅将其作为候选思路，
最终是否采用由参赛人员结合赛题条件、数学合理性和实际计算结果判断。

对于AI生成的代码，参赛人员重点进行了以下人工检查和修改：

\begin{itemize}
    \item 检查10分钟时间索引及附件模板中的时间标签对应关系；
    \item 检查功率由kW转换为时段电量kWh时是否正确乘以$1/6$小时；
    \item 检查储能状态递推、容量范围和最大充放电功率约束；
    \item 检查计划购电、紧急购电以及调整购电的费用倍数和结算方式；
    \item 检查问题二至问题四的决策过程中是否错误使用尚未获得的未来实际数据；
    \item 检查结果工作簿中的日期、时段、电量和费用是否与程序计算结果一致。
\end{itemize}

经人工检查后再运行程序，并利用实际计算结果进一步核对供需平衡、储电量连续性、
上下界、充放电功率以及费用分项。问题一的购电费用和全天购电量还进行了独立复算，
得到购电费用$35126.9486$元、全天购电量$59482.6990$ kWh，与程序结果一致。

论文中的数值结果均来自最终实际运行的程序输出，而不是直接采用AI生成的数值。
AI生成的文字、公式说明和图表说明在进入论文前均根据实际模型和运行结果进行了修改，
删除了与最终方法不一致、无法由计算结果支持或超出赛题条件的内容。

因此，本次AI的主要作用是辅助理解题意、将人工确定的模型转化为代码、协助检查程序
及整理表达；模型约束的确定、程序核验、实际计算、策略选择和最终结果确认仍由参赛人员完成。

\end{document}